\section{МОДЕЛЬ ЭРДЁША--РЕНЬИ. ЕЁ СВОЙСТВА}
\label{g1}
В данном разделе вводятся ключевые определения и обозначения, 
излагается алгоритм поиска в глубину, используемый для нахождения 
компонент связности графа, и формулируются теоретические результаты
о свойствах модели $G(n,p)$.
\subsection{Основные определения}

Введем понятие неориентированного графа: $G=(V,E)$, где  $V = \{1, 2, \dots, n\}$ множество натуральных чисел от 
$1$ до  $n$, $E$ -- множество неориентированных ребер. Определим
путь как последовательность вершин графа, такую, что
две соседние вершины связаны ребром и каждое ребро встречается не более одного раза. Циклом называется такой путь, 
который начинается и заканчивается в одной вершине. Граф является связным, если между любыми двумя вершинами существует 
путь. Связный граф без циклов называется деревом. Подграф графа $G = (V,E)$ -- это такой граф $G' =(V',E')$,
$V' \subset V$, $E' \subset E$. Компонентой связности графа 
называется максимальный по включению связный подграф, добавление вершин в который 
нарушит связность.
\subsection{Поиск компонент связности с помощью поиска в глубину}
Рассмотрим алгоритм поиска компонент связности с помощью поиска в 
глубину. Заданы три множества $S$ -- множество посещенных вершин,
 $T$ -- множество непосещенных вершин и упорядоченное множество $U = V \setminus (S \cup T)$, организованное как стек. В начале работы алгоритма $T=V$,
 $S= U = \emptyset$. Алгоритм прекращает свою работу в случае, если
 $U \cup T = \emptyset$. Если $U$ не пусто, $v$-- последняя добавленная в  $U$ вершина, то в  $T$ 
 берется смежная вершина с $v$ вершина, удаляется из  $T$,   
 добавляется в  $U$. Если такой вершины нет, то  $v$ удаляется 
 из  $U$ и добавляется в  $S$. Если  $U$ пусто, то из 
 $T$ удаляется первая вершина и добавляется в  $U$. 
 Компонента связности находится тогда, когда множество  $U$
 становится пустым после удаления вершины.
\subsection{Свойства модели Эрдёша--Реньи}

Рассмотрим модель случайного графа  $G(n,p)$, $n$ -- количество 
вершин,  $p \in [0;1]$-- вероятность, с которой каждое ребро добавляется независимо друг от друга в граф. 
Будем говорить, что событие $\mathcal{E}_{n}$ 
происходит с высокой вероятностью, если $\lim_{n \to \infty} \text{Pr}(G \sim G(n,p) \in \mathcal{E}_{n}) =1 $. Иными
словами, вероятность события стремится к $1$ при $n \to \infty$. В статье Кривелевича и Судакова \cite{ft} были доказаны оценки размеров 
компонент связности и длины пути. Пусть $\epsilon > 0$ -- 
достаточно малое число, тогда если граф  $G$ 
создан в соответствии с моделью  $G\left(n,\frac{1-\epsilon}{n}\right)$,
то с высокой вероятностью размер всех его компонент связности не превосходит 
$\frac{7}{\epsilon^2} \ln{n}$. Если граф $G$ 
создан в соответствии с моделью  $G\left(n,\frac{1+\epsilon}{n}\right)$, то с высокой вероятностью 
этот граф содержит путь длиной хотя бы $\frac{\epsilon^2 n}{5}$ 
и имеет компоненту связности, размер которой больше чем 
$\frac{en}{2}$. Значение $p=\frac{1}{n}$ является критическим , при его достижения
структура графа меняется. При $p< \frac{1}{n}$ компоненты 
связности имеют небольшой размер, среди них много деревьев, 
после преодоления критического значения граф приобретает гигантскую компоненту связности. Также в книге Райгородского
\cite{rai} доказаны следующие свойства, если граф создан 
в соответствии с моделью $G\left(n,c\frac{\ln{n}}{n}\right)$,
то граф с высокой вероятностью связен, если $c > 1$ и несвязен, если $c < 1$. Значение $p = \frac{\ln{n}}{n}$ является фазовым переходом, после его достижение структура графа меняется. Таким 
образом в модели $G(n,p)$ существуют два 
критических значения  $p$, при которых происходит фазовый переход.
После достижения  $p = \frac{1}{n}$ появляется гигантская компонента связности, после достижения $p = \frac{\ln{n}}{n}$ 
граф становится связным.
